\section{Projection areas and longest chords of simplices}
\label{sec:projections}

The main constructions use the direct projection calculation in
Section~\ref{sec:tetrahedron}. This appendix gives a more general
normalization and two alternative geometric explanations of the regular
tetrahedron's minimum shadow. They also fix the Euclidean metric for the
coordinate changes in Appendix~\ref{sec:coordinate-realizations}.

The simplex chord--projection identity below is attributed to Martini and
Weissbach \cite{MartiniWeissbach1992} by Heinrich, who states the exact
constant \cite[p.~8268]{Heinrich2014}. Martini characterizes simplices
among convex polytopes in dimension at least three by the independence
of the chord--projection product from direction
\cite[Theorem~2(C)]{Martini1991}. We compute the constant directly.

\begin{proposition}[Simplex chord--projection identity]
\label{prop:simplex}
Let $S\subset\R^d$, $d\ge2$, be a full-dimensional simplex, and let
$u\in\R^d$ be a unit vector. Write $V$ for its $d$-dimensional volume,
$A(u)$ for the $(d-1)$-dimensional volume of its orthogonal projection
onto $u^\perp$, and $L(u)$ for the greatest length of a chord of $S$
parallel to $u$. Then
\[
 A(u)L(u)=dV,\qquad
 \min_{\norm{u}=1}A(u)=\frac{dV}{\operatorname{diam}S}.
\]
\end{proposition}
\begin{proof}
Let $\lambda_0,\ldots,\lambda_d$ be the affine barycentric coordinates
of $S$, and put $b_i=\partial_u\lambda_i$. They satisfy
$\sum_i b_i=0$ and are not all zero, because the barycentric coordinate
map is injective. Hence
\[
 c=\sum_i\max(0,-b_i)=\sum_i\max(0,b_i)>0.
\]
For $s\ge0$, both $x$ and $x+su$ belong to $S$ exactly when
\[
 \lambda_i(x)\ge s\max(0,-b_i)\quad(0\le i\le d),
 \qquad \sum_i\lambda_i(x)=1.
\]
For $0\le s<1/c$, this intersection is a translate of $(1-sc)S$;
at $s=1/c$ it is one point, and above that value it is empty.
In particular, $L(u)=1/c$.

A fiber of $S$ parallel to $u$ contains two points separated by $su$
if and only if its length is at least $s$. Thus, for $0<s<L(u)$,
the projection of $S\cap(S-su)$ is the superlevel set of the fiber
length at $s$, and its $(d-1)$-volume is
$A(u)(1-s/L(u))^{d-1}$. Integrating the fiber lengths gives
\[
 V=A(u)\int_0^{L(u)}(1-s/L(u))^{d-1}\,ds
   =\frac{A(u)L(u)}{d}.
\]
This is the layer-cake formula, obtained directly by integrating
$\ell=\int_0^\infty\mathbf1_{\{s<\ell\}}\,ds$ on each fiber.
Finally, the largest chord length over all directions is the diameter:
every chord joins two points of $S$, and a pair of points realizing
the diameter is itself a chord. Taking the minimum of $dV/L(u)$
proves the second formula.
\end{proof}

\label{par:entering-faces}
There is a short alternative proof of the same constant using the
faces through which the fibers enter. If the facet opposite vertex $i$
has $(d-1)$-volume $S_i$ and altitude $h_i$, then
$V=S_i h_i/d$ and $\norm{\nabla\lambda_i}=1/h_i=S_i/(dV)$.
The facets with $b_i>0$ are the entering facets for direction $u$.
Their projections partition the shadow apart from a null set, and
facet $i$ contributes
\[
 S_i\frac{b_i}{\norm{\nabla\lambda_i}}=dVb_i.
\]
Consequently $A(u)=dV\sum_{b_i>0}b_i=dVc=dV/L(u)$.
This proof also includes directions parallel to a facet, whose
contribution is zero.

For any simplex, the distance between two convex combinations of its
vertices is at most the largest distance between two vertices. Hence the
diameter of a regular tetrahedron is its edge length. For edge length
$\ell>0$, the volume is
$\ell^3/(6\sqrt2)$ and the diameter is $\ell$. Therefore
\begin{equation}\label{eq:regular-normalization}
 A_{\min}=\frac{\ell^2}{2\sqrt2},\qquad
 \frac{A_{\min}}2=\frac{\ell^2}{4\sqrt2}.
\end{equation}
For example, the volume follows by taking a face of area
$\sqrt3\ell^2/4$ and altitude $\sqrt{2/3}\ell$.
A similarity of ratio $r>0$ multiplies every perpendicular base area
and $A_{\min}$ by $r^2$, while preserving all cylinder counts and
both normalized covering costs.

\subsection{A different reason that an edge direction minimizes the shadow}

The chord identity is not needed to locate the minimum for a regular
tetrahedron. There is an independent argument based on the piecewise
linear expression in Cauchy's face formula~\eqref{eq:cauchy}.
Indeed, some minimizing direction for any tetrahedron is parallel to an
edge; for a regular tetrahedron all edge directions give the same area.
To see this directly, argue as follows.
The continuous positive projection-area function $A$ attains its
minimum on the unit sphere. Suppose that a minimizing $u$ is
perpendicular to at most one facet normal. Choose a unit $v\perp u$,
also perpendicular to that normal if it exists. For sufficiently small
$\theta>0$, all nonzero signs in the face formula stay unchanged at
\[
 u_\pm=(\cos\theta)u\pm(\sin\theta)v,
\]
and a zero scalar product remains zero. Consequently
\[
 A(u_+)+A(u_-)=2\cos\theta\,A(u)<2A(u),
\]
contradicting minimality. Equivalently, along this great-circle arc the
area is a fixed linear form and satisfies $A''=-A<0$, so it cannot
have a local minimum. A minimizing direction is therefore perpendicular
to at least two facet normals and hence parallel to the common edge.
For a regular tetrahedron, isometries permute the edges. In the edge-$2$
model
\[
 \{(x,y,\sqrt2 t):0\le t\le1,\ |x|\le t,\ |y|\le1-t\},
\]
projection in the $x$-direction has area
$\sqrt2\int_0^1 2(1-t)\,dt=\sqrt2$. Scaling gives
$\ell^2/(2\sqrt2)$ for every edge direction.

\paragraph{The regular normal frame.}
A further algebraic check, useful in barycentric coordinates, comes
from the regular tetrahedron's normal frame.\label{par:normal-frame}
If its outward unit
normals are $n_1,\ldots,n_4$, then
\[
 \sum_i n_i=0,\qquad \sum_i n_i n_i^{\mathsf T}=\frac43 I.
\]
These identities follow, for example, from the four vectors
$(1,1,1),(1,-1,-1),(-1,1,-1),(-1,-1,1)$ after normalization.
For a unit direction $u$, set $s_i=n_i\cdot u$ and let $P$ be the
sum of their positive values. The negative values have absolute sum
$P$, so
\[
 \frac43=\sum_i s_i^2\le2P^2.
\]
Since a facet has area $\sqrt3\ell^2/4$, the positive-face form of
Cauchy's formula gives
$A(u)=(\sqrt3\ell^2/4)P\ge\ell^2/(2\sqrt2)$.
In the barycentric realization\label{par:barycentric-projection}
$\{(a,b,c,d)\in[0,\infty)^4:a+b+c+d=1\}$, the Euclidean
metric is the one induced from $\R^4$, the edge length is $\sqrt2$,
and the same calculation reads
$A(u)=\tfrac12\sum_i|u_i|\ge1/\sqrt2$ for unit sum-zero vectors $u$.
These equivalent frame calculations fix the metric explicitly when
we change coordinates below.
